%Chargement des paquets
\usepackage{amsmath}
\usepackage{amsfonts}
\usepackage{amssymb}
\usepackage{enumerate}
\usepackage{mathtools}
\usepackage{bbm}
\usepackage{xparse, etoolbox}
\usepackage{enumerate}
\usepackage{enumitem}
\usepackage{mathabx}
\usepackage{babel}[french]

%environment
\usepackage{amsthm}
\usepackage{keytheorems}
\usepackage{hyperref} 

%Interface théorème
\newkeytheorem{lem}[
	name = Lemme
]

\newkeytheorem{prop}[
	name = Proposition
]

\newkeytheorem{thm}[
	name = Théorème
]

\newkeytheorem{coro}[
	name = Corollaire
]

\renewcommand*{\proofname}{Démonstration}

\theoremstyle{definition}
\newtheorem{defn}{Définition}

\theoremstyle{definition}
\newtheorem*{nota}{Notation}

\theoremstyle{definition}
\newtheorem*{limi}{Liminaire}

\theoremstyle{remark}
\newtheorem{rmq}{Remarque}

\theoremstyle{plain}
\newtheorem*{fait}{Fait}


%Ensembles de nombres
\newcommand{\N} {\mathbb{N}}
\newcommand{\Ne}{\N^\ast}
\newcommand{\Z} {\mathbb{Z}}
\newcommand{\Q} {\mathbb{Q}}
\newcommand{\R} {\mathbb{R}}
\newcommand{\Rb}{\overline{\mathbb{R}}}
\newcommand{\Rp}{\R_+}
\newcommand{\Rm}{\R_-}
\newcommand{\K} {\mathbb{K}}
\newcommand{\Ket} {\mathbb{K^{\ast}}}
\newcommand{\Cx}{\mathbb{C}}

%Ensembles, tribus
\newcommand{\A}{\mathbb{A}}
\renewcommand{\P}{\mathcal{P}}
\newcommand{\T}{\mathcal{T}}
\newcommand{\F}{\mathcal{F}}
\newcommand{\C} {\mathcal{C}}
\newcommand{\ouv}{\mathcal{O}}
\newcommand{\B}{\mathcal{B}}
\newcommand{\M}{\mathcal{M}}
\newcommand{\etag}{\mathcal{E}}
\newcommand{\etagOp}{\etag^{+}(\Omega)}
\newcommand{\etagO}{\etag(\Omega)}
%%Lp avant quotientage
\newcommand{\Lic}{\mathcal{L}^1}
\newcommand{\Lpc}{\mathcal{L}^p}
\newcommand{\Linfc}{\mathcal{L}^{\infty}}
\newcommand{\Lifull}{\Lic(\Omega, \B, m)}
\newcommand{\Lpfull}{\Lpc(\Omega, \B, m)}
\newcommand{\Linffull}{\Linfc(\Omega, \B, m)}
%%Lp après quotientage
\newcommand{\Li}{L^1}
\newcommand{\Ld}{L^2}
\newcommand{\Lp}{L^p}
\newcommand{\Linf}{L^{\infty}}
\newcommand{\Listd}{\Li(\Omega, m)} 
\newcommand{\Ldstd}{\Ld(\Omega, m)} 
\newcommand{\Lpstd}{\Lp(\Omega, m)} 

%Quelques opérateurs
\DeclareMathOperator{\codim}{codim}
\DeclareMathOperator{\rg}{rg}
\DeclareMathOperator{\tr}{tr}
\DeclareMathOperator{\vect}{Vect}
\DeclareMathOperator{\ima}{Im}
\DeclareMathOperator{\id}{Id}
\DeclareMathOperator{\sign}{sign}
\DeclareMathOperator{\ext}{ext}
\newcommand{\ind}{\mathbbm{1}}
\newcommand*{\spr}[2]{\left(\,#1\,\middle|\,#2\,\right)}
\newcommand{\interior}[1]{%
  {\kern0pt#1}^{\mathrm{o}}%
}
\DeclareMathOperator{\card}{card}
\DeclareMathOperator{\ppcm}{ppcm}

%Commandes de pure flemme
\newcommand{\veps}{\varepsilon}
\newcommand{\intab}{\int_{a}^{b}}
\newcommand{\intR}{\int_{-\infty}^{\infty}}
\newcommand{\intinf}{\int_{- \infty}^{+ \infty}}
\newcommand{\equi}{\Leftrightarrow}
\newcommand{\Kev}{$\K$-espace vectoriel }
\newcommand{\Kevs}{$\K$-espaces vectoriels }
\newcommand{\Rev}{$\R$-espace vectoriel }
\newcommand{\Revs}{$\R$-espaces vectoriels }
\newcommand{\Cev}{$\Cx$-espace vectoriel }
\newcommand{\Cevs}{$\Cx$-espaces vectoriels }
\newcommand{\evn}{(E, \norm{.})}
\newcommand{\interf}[2]{[#1,#2]}
\newcommand{\limb}{\lim\limits_}
\newcommand{\lebn}{\lambda_n}
\newcommand{\pp}{\text{~presque partout}}
\newcommand{\derivp}[2]{\frac{\partial #1}{\partial #2}}
\newcommand{\famiu}{(u_i)_{i \in I}}
\newcommand{\sev}{\text{~s.e.v.}}
\newcommand{\Mnp}{M_{n,p}(\K)}
\newcommand{\Mn}{M_{n}(\K)}
\newcommand{\tq}{\text{~t.q.~}}
\newcommand{\mq}{~montrer~que~}
\newcommand{\et}{\quad \text{et} \quad}
\newcommand{\ou}{\text{~ou~}}
\newcommand{\Kx}{\K[X]}


%Commandes de gars chelou
\renewcommand{\subset}{\subseteq}

%Commande restriction, ty duckduckgo
\def\restr#1#2{\mathchoice
              {\setbox1\hbox{${\displaystyle #1}_{\scriptstyle #2}$}
              \restraux{#1}{#2}}
              {\setbox1\hbox{${\textstyle #1}_{\scriptstyle #2}$}
              \restraux{#1}{#2}}
              {\setbox1\hbox{${\scriptstyle #1}_{\scriptscriptstyle #2}$}
              \restraux{#1}{#2}}
              {\setbox1\hbox{${\scriptscriptstyle #1}_{\scriptscriptstyle #2}$}
              \restraux{#1}{#2}}}
\def\restraux#1#2{{#1\,\smash{\vrule height .8\ht1 depth .85\dp1}}_{\,#2}} 

%Quelques normes
\DeclarePairedDelimiter\abs{\lvert}{\rvert}
\DeclarePairedDelimiter\labs{\bigl\lvert}{\bigr\rvert}
\DeclarePairedDelimiter\norm{\lVert}{\rVert}
\DeclarePairedDelimiterXPP\normi[1]{}\lVert\rVert{_1}{\ifblank{#1}{\:\cdot\:}{#1}}
\DeclarePairedDelimiterXPP\normd[1]{}\lVert\rVert{_2}{\ifblank{#1}{\:\cdot\:}{#1}}
\DeclarePairedDelimiterXPP\normp[1]{}\lVert\rVert{_p}{\ifblank{#1}{\:\cdot\:}{#1}}
\DeclarePairedDelimiterXPP\normq[1]{}\lVert\rVert{_q}{\ifblank{#1}{\:\cdot\:}{#1}}
\DeclarePairedDelimiterXPP\norminf[1]{}\lVert\rVert{_\infty}{\ifblank{#1}{\:\cdot\:}{#1}}

%Bracket en angle
\DeclarePairedDelimiter\angl{\langle}{\rangle}

%Set, parenthèses
\newcommand{\set}[1]{\{ #1 \}}
\newcommand{\bigset}[1]{\bigl\{ #1 \bigr\}}
\newcommand{\Bigset}[1]{\Bigl\{ #1 \Bigr\}}
\newcommand{\bigpar}[1]{\bigl( #1 \bigr)}
\newcommand{\Bigpar}[1]{\Bigl( #1 \Bigr)}

%Opitmisation des opérateurs de comparaison
\DeclarePairedDelimiter\tio{o (}{)}
\DeclarePairedDelimiter\gro{O (}{)}


\pagestyle{fancy}

\fancyhead[C]{}
\fancyhead[R]{}

\begin{document}

\underline{Source :} Grifone - Algèbre linéaire - page 252 puis page 286 

\section{Cas euclidien}

Dans cette section $E$ est un espace euclidien (de dimension $n \geq 1$) muni d'un produit scalaire dénoté par $\angl{ \cdot, \cdot}$.

\begin{defn}[Endomorphisme autoadjoint]
Un endomorphisme $f$ de $E$ est dit autoadjoint si :
\[ \forall x,y \in E, \quad \angl{f(x), y} = \angl{x, f(y)} . \]
\end{defn}

\begin{lem}
\label{lem:1}
Un endomorphisme d'un espace euclidien est autoadjoint si, et seulement si, sa matrice dans une base orthonormée est symétrique.
\end{lem}

\begin{proof}
Soit $f$ un endomorphisme autoadjoint, notons $A$ sa matrice dans une base orthonormée de $E$. 
Dans cette même base, on a :
\[ \forall X, Y \in M_{n,1}(\R), \quad (AX)^T Y = X^T (AY) , \]
soit :
\[ \forall X, Y \in M_{n,1}(\R), \quad X^TA^T Y = X^T AY , \]
ce qui est équivalent à $A = A^T$. 
La réciproque est évidente. 
\end{proof}

\begin{thm}[Théorème spectral euclidien]
Soit $f$ un endomorphisme autoadjoint d'un espace euclidien. Alors :
\begin{enumerate}[label=(\roman*)]
\item $f$ est diagonalisable ;
\item les sous espaces propres de $f$ sont deux à deux orthogonaux. En particulier, on peut construire une base orthonormée 
de vecteurs propres en choisissant une base orthonormée dans chaque espace propre.
\end{enumerate}
\end{thm}

\begin{proof}
\begin{enumerate}[label=(\alph*)]

\item
On muni $E$ d'une base orthonormée, et on note $A$ la matrice de $f$ dans cette base. 
On montre d'abord que le polynôme caractéristique $P_f$ de $f$ est scindé, c'est-à-dire que toutes ses valeurs propres sont réelles.
Le corps $\Cx$ étant algébriquement clos, $P_f$ admet au moins une racine complexe que l'on note $\lambda$.
On veut montrer que $\lambda \in \R$. 
Soit $x$ un vecteur propre associé à $\lambda$ et $X$ sa matrice colonne (à coefficients éventuellements complexes) dans la base
choisie. Comme $A$ est symétrique, on a :
\[ (AX)^T \widebar{X} = X^T A \widebar{X} . \]
D'autre part, $AX = \lambda X$ donc $\widebar{AX} = \widebar{A} \widebar{X} = \widebar{\lambda} \widebar{X}$, 
et, comme $A$ est réelle :
\[ (\lambda X)^T \widebar{X} = X^T \widebar{\lambda} \widebar{X}, \]
autrement dit :
\[ \lambda (\abs{x_1}^2 + \dots + \abs{x_n}^2) = \widebar{\lambda} (\abs{x_1}^2 + \dots + \abs{x_n}^2) , \]
et donc $\lambda =  \widebar{\lambda}$ puisque $v \neq 0$. Les valeurs propres de $f$ sont donc toutes réelles 
(et par conséquence, les vecteurs propres sont à coefficients réels dans la base orthonormée).

\item
On va montrer, par récurrence sur la dimension $n$ de $E$, qu'il existe une base de $E$ formées de vecteurs propres de $f$. \\
Pour $n=1$, la propriété est trivialement vraie.
Supposons la vraie pour $n-1 \geq 1$ quelconque et montrons qu'elle reste vraie au rang $n$. 
Soit $\lambda$ une valeur propre de $f$ et $v$ un vecteur propre correspondant. 
On pose $H = \vect(v)^{\perp}$, on remarque que $\dim{H} = n-1$ (car $H$ est le noyau de la forme linéaire $x \mapsto\angl{x,v}$). \\
De plus, $H$ est stable par $f$, en effet si $x \in H$, alors :
\[ \angl{f(x), v} = \angl{x, f(v)} = \lambda \angl{x, v} = 0 . \]
Ainsi, la restriction de $f$ à $H$ est un endomorphisme de $H$. \\
Par ailleurs, $\restr{f}{H}$ est un endomorphisme autoadjoint de $H$ (pour la restriction du produit scalaire à $H$), en effet :
\[ \angl{\restr{f}{H}(x), y}_H = \angl{f(x) , y} = \angl{x, f(y)} = \angl{x, \restr{f}{H}(y)}_H \quad \forall x, y \in H  ,\]
(il s'agit d'un simple jeu d'écriture). Or, $\dim{H} = n-1$, donc, d'après l'hypothèse de récurrence, il existe une base $\B$ formée
des vecteurs propres de $\restr{f}{H}$. Comme $E = H \oplus \set{x}$, 
la base $\B \cup \set{x}$ est bien une base de $E$ formée des vecteurs propres de $f$. 
Ainsi, la propriété est héréditaire, donc vraie pour tout $n \geq 1$. \\
Ce qui permet d'affirmer que $f$ est diagonalisable.

\item
Enfin, on va démontrer que les sous-espaces propres de $f$ sont deux à deux orthogonaux. 
Soient $\lambda_1$ et $\lambda_2$ deux valeurs propres distinctes de $f$ et $v_1, v_2$ les vecteurs propres correspondants.
Il s'agit de montrer que $\angl{v_1, v_2} = 0$. On a :
\[ \angl{f(v_1), v_2} = \lambda_1 \angl{v_1, v_2} \et \angl{v_1, f(v_2)} = \lambda_2 \angl{v_1, v_2} . \]
Or $f$ est autoadjoint, donc $(\lambda_1 - \lambda_2) \angl{v_1 , v_2} = 0$ 
soit $\angl{v_1 , v_2} = 0$, puisque $\lambda_1 \neq \lambda_2$.

\end{enumerate}
\end{proof}

D'après le lemme $\ref{lem:1}$ on peut donner une formulation matricielle de ce théorème :

\begin{thm}
Toute matrice symétrique réelle est diagonalisable dans $\R$ et ses sous-espaces propres sont deux à deux orthogonaux pour le 
produit scalaire canonique de $\R^n$.
\end{thm}

De plus, comme une matrice de passage d'une base orthonormée à une autre est une matrice orthogonale, on peut (encore) 
reformuler ce théorème :

\begin{thm}
Si $A$ est une matrice symétrique réelle, alors il existe une matrice orthogonale $P$ telle que la matrice $P^T A P$ est diagonale.
\end{thm}

\section{Cas hermitien}

\begin{nota}
Dans cette section $E$ est un espace hermitien (de dimension $n \geq 1$) muni d'un produit scalaire dénoté par $\angl{.,.}$. \\
Pour $f$ un endomorphisme de $E$, $f^{\ast}$ désigne l'endomorphisme adjoint de $f$. \\
Pour $A \in M_n(\C), A^{\ast}$ désigne la matrice transconjugée de $A$. \\
Dans un espace hermitien, l'analogue des matrices othogonales sont les matrices unitaires.
\end{nota}

\begin{defn}[Endomorphisme normal]
Un endomorphisme $f$ d'un espace hermitien est dit normal si :
\[ f^{\ast} \circ f = f \circ f^{\ast} . \]
\end{defn}

\begin{rmq}
De façon équivalente, une matrice $A \in M_n(\C)$ est dite normale si :
\[ A^{\ast} A = A A^{\ast} . \]
Les matrices normales sont les matrices qui représentent les endomorphismes normaux dans une base orthonormée.
\end{rmq}

\begin{lem}
\label{lem:2}
Un endomorphisme $f$ est normal si, et seulement si :
\[ \forall x, y \in E, \quad \angl{ f(x), f(y) } = \angl{ f^{\ast}(x), f^{\ast}(y) } . \] 
\end{lem}

\begin{proof}
Simple jeu d'écriture, par exemple pour la condition nécessaire, soient $x, y \in E$.
\begin{align*}
\angl{ f(x), f(y) } & = \angl{ x, f^{\ast} \circ f(y) } \\
& = \angl{ x, f \circ f^{\ast}(y) } \\
& = \angl{ f^{\ast}(x),  f^{\ast}(y) } .
\end{align*}
\end{proof}

\begin{thm}
Soit $f$ un endomorphisme normal d'un espace hermitien. Alors :
\begin{enumerate}[label=(\roman*)]
\item $f$ est diagonalisable ;
\item les sous espaces propres de $f$ sont deux à deux orthogonaux. En particulier, on peut construire une base orthonormée 
de vecteurs propres en choisissant une base orthonormée dans chaque espace propre.
\end{enumerate}
\end{thm}

\begin{proof}
On démontre ce résultat par récurrence sur la dimension $n$ de $E$. \\
Il est clairement vrai pour $n=1$, on le suppose vrai jusqu'à l'ordre $n-1 \geq 1$ et on démontre qu'il reste vrai à l'ordre $n$.
Soit $\lambda$ une valeur propre de $f$ et $E_{\lambda}$ le sous-espace propre corresondant, on peut écrire :
\[ E = E_{\lambda} \oplus E_{\lambda}^{\perp} , \]
avec $\dim{E_{\lambda}^{\perp}} \geq n-1$. On procède comme dans le cas euclidien.

\begin{enumerate}[label=(\alph*)]

\item
On veut montrer que $E_{\lambda}^{\perp}$ est stable par $f$. On note d'abord que $E_{\lambda}$ est stable par $f^{\ast}$.
En effet, pour $x \in E_{\lambda}$, on a :
\[ f \circ f^{\ast} (x) = f^{\ast} \circ f (x) = \lambda f^{\ast} (x) . \]
Prenons maintenant $y \in E_{\lambda}^{\perp}$, alors pour tout $x \in E_{\lambda}$ :
\[ \angl{ f(y) , x } = \angl{ y , f^{\ast}(x) } = 0 , \]
par stabilité de $E_{\lambda}$ par $f^{\ast}$. Ce qui revient à dire que $f(y) \in E_{\lambda}^{\perp}$.

\item
On veut ensuite montrer que la restriction de $f$ à $E_{\lambda}^{\perp}$ est un endomorphisme normal.
On note d'abord que $E_{\lambda}^{\perp}$ est stable par $f^{\ast}$. 
En effet, pour $y \in E_{\lambda}^{\perp}$ fixé et pour tout  $x \in E_{\lambda}$ on a :
\[ \angl{ y , f(x) } = \lambda \angl{ y, x} = 0 = \angl{ f^{\ast}(y), x} . \]
En utilisant le lemme $\ref{lem:2}$ et la stabilité de $E_{\lambda}^{\perp}$ par $f$ et $f^{\ast}$ on peut écrire 
$\forall x, y \in E_{\lambda}^{\perp}$ :
\begin{align*}
\angl{ \restr{f}{E_{\lambda}^{\perp}}(x) , \restr{f}{E_{\lambda}^{\perp}}(y)}_{E_{\lambda}^{\perp}} & = \angl{ f(x), f(y) } \\
& = \angl{ f^{\ast}(x), f^{\ast}(y) } \\
& =  \angl{ \restr{f}{E_{\lambda}^{\perp}}^{\ast}(x) , \restr{f}{E_{\lambda}^{\perp}}^{\ast}(y)}_{E_{\lambda}^{\perp}} ,
\end{align*}
autrement dit, toujours d'après le lemme $\ref{lem:2}$, $\restr{f}{E_{\lambda}^{\perp}}$ est normale.

\item 
On peut maintenant conclure. La propriété est vraie (par hypothèse) pour $\restr{f}{E_{\lambda}^{\perp}}$, qui est donc diagonalisable
et dont les sous-espaces propres sont deux à deux orthogonaux. 
$E = E_{\lambda} \oplus E_{\lambda}^{\perp}$ la propriété est clairement vraie pour $f$, ce qui achève la récurrence.

\end{enumerate}
\end{proof}

Un endomorphisme autoadjoint étant clairement normal, on en tire le corollaire suivant.

\begin{coro}[Théorème spectral hermitien]
Soit $f$ un endomorphisme autoadjoint d'un espace hermitien. Alors :
\begin{enumerate}[label=(\roman*)]
\item les valeurs propres de $f$ sont toutes réelles ;
\item $f$ est diagonalisable ;
\item les sous espaces propres de $f$ sont deux à deux orthogonaux. En particulier, on peut construire une base orthonormée 
de vecteurs propres en choisissant une base orthonormée dans chaque espace propre.
\end{enumerate}
\end{coro}

\begin{proof}
Seul le point (i) est à démontrer (les deux autres sont conséquences du théorème précédent). Soit $\lambda$ une valeur propre de $f$
et $v$ un vecteur propre associé, on a :
\[ \angl{ f(v) ,v } = \widebar{\lambda} \norm{v} \et \angl{ f(v) ,v } = \angl{v, f(v)} = \lambda \norm{v} . \]
Do'ù $\lambda \in \R$ (car $v$ non nul).
\end{proof}

Ou encore, en s'appuyant sur le lemme ci-dessous (analogue au lemme $\ref{lem:1}$ du cas euclidien les formulations matricielles.

\begin{lem}
Un endomorphisme d'un espace hermitien est autoadjoint si, et seulement si, sa matrice dans une base orthonormée est hermitienne.
\end{lem}

\begin{thm}
Toute matrice hermitienne est diagonalisable dans $\R$ et ses sous-espaces propres sont deux à deux orthogonaux pour le 
produit scalaire hermitien canonique de $\C^n$.
\end{thm}

\begin{thm}
Si $A$ est une matrice hermitienne, alors il existe une matrice unitaire $U$ telle que la matrice $U^{\ast} A U$ est diagonale réelle.
\end{thm}

\begin{rmq}
Dans un espace hermitien, les endomorphismes autoadjoints sont aussi appelés... endomorphismes hermitiens. 
\end{rmq}

\end{document}